\section{Introduction}
\label{sec:intro}

A backdoored classifier behaves normally on clean inputs and predicts an attacker's class whenever a trigger is present in the input. The attacker plants the behaviour by poisoning a small share of the training data, often 1\% or less, and the defender receives the trained model without the training set. Input-level detection asks the most practical question in this setting: given the deployed model and 1 incoming input, is that input carrying a trigger? A detector that answers well can run in front of the model in production, needs no retraining, and can be combined with removal defences later.

Prediction Shift Backdoor Detection (PSBD) \cite{li2025psbd} is one such detector. Its observation is that a backdoored input is decided by a shortcut the model learned quickly and firmly, while a clean input is decided by a broad set of features. When dropout is switched on at inference, the broad decision wobbles and the shortcut does not. PSBD runs the model a few times with dropout active, measures how much the probability of the predicted class falls, and flags inputs whose prediction barely moves. It needs only a small clean validation set to set the threshold. On ResNet-18 and VGG-16 the method reports strong results across 7 attacks.

Vision Transformers now sit in most vision pipelines, and the PSBD paper lists them as future work. Nothing about the prediction shift argument is specific to convolutions, but 2 things change on a transformer. First, dropout can be placed at many distinct sites inside a block, and these sites are not interchangeable: the input to an attention block, the input to its MLP, the branch output before the residual add and the residual stream after the add each see a different tensor with a different role. On a ResNet basic block there is essentially 1 natural site, and the original paper used it. Second, dropout is not the only way to perturb a transformer. Whole tokens can be masked, channels can be masked, noise can be added to a LayerNorm's input or its gain can be scaled, and each of these interacts differently with the attention and normalisation that follow. So adapting PSBD to a ViT is a search over 2 axes, the position and the operator, rather than a port.

We ran that search. On 65 backdoored ViT-B/16 checkpoints covering 4 datasets, 7 attacks and 3 poison rates, we swept 18 placements, each a position paired with an operator, at a ladder of perturbation rates, and read every placement on the same held-out split with the same threshold rule. The result is a placement rather than an operator: masking whole tokens at the input of each attention block detects with a mean AUROC of 0.935 against 0.832 for the literal port of the original placement, and the gain concentrates where detection is hardest. Knowing which placement wins is useful. Knowing why it wins is what lets the finding transfer, so the second half of the paper follows the backdoor through the network with activation patching, direction ablation and a per-layer read of the class token, and arrives at an account that also predicts which placements lose, which attacks show the prediction shift phenomenon the original paper describes, and how the picture changes on Swin.

The contributions are as follows.
\begin{enumerate}
  \item A systematic adaptation of PSBD to ViT-B/16 over a fixed basis of 18 placements, 65 checkpoints and a single decision rule, showing that token masking before the attention LayerNorm gains +0.103 mean AUROC over the published placement, +0.110 on hard attacks and +0.161 at 1\% poisoning (\cref{sec:results}).
  \item A separation of position from operator, showing that both matter and neither reduces to a family: the same operator moves by 0.109 between the attention and the MLP input, and the same site moves by 0.206 between masking and noise, with the sign of the operator effect reversing between the 2 sites (\cref{sec:results:operator}).
  \item A mechanistic account, tested on the same checkpoints, in which the backdoor is 1 direction routed through attention from the trigger's tokens to the class token in the last third of the network, which predicts the winning placement, the operator reversal through LayerNorm absorption, and the attacks on which the original paper's prediction shift phenomenon does and does not appear (\cref{sec:mechanism}).
  \item A transfer of the placement to Swin-S with a gain of +0.091 over the published placement, together with the negative and qualified results a deployer needs: the adaptive attacker who trains against 1 probe defeats it and only a union over probes restores detection, forward passes beyond 3 help mainly at 1\% poisoning, and all-to-all attacks remain out of reach (\cref{sec:swin,sec:robustness,sec:limitations}).
\end{enumerate}
Every number in this paper is produced by a script from a coverage ledger that declares the panel, and the code, the ledger and the generators are released.
